% \iffalse meta-comment
%
% Copyright 2025-2026 Zhenyu Zhong
%
% This work may be distributed and/or modified under the
% conditions of the LaTeX Project Public License, either version 1.3
% of this license or (at your option) any later version.
% The latest version of this license is in
%   https://www.latex-project.org/lppl.txt
% and version 1.3c or later is part of all distributions of LaTeX
% version 2008 or later.
%
% \fi
%
% \subsubsection{中文语言数据}
%
% 中文是论文的默认语言，除模式二标题外的设置均与 \cls{ctexbook} 的默认值一致。
%
% \begin{macro}{\@@_lang_zh_names:}
%
%    \begin{macrocode}
\cs_new_protected:Nn \@@_lang_zh_names:
    {
        \cs_set:Npn \prefacename         { 前言 }
        \cs_set:Npn \symabbrname         { 符号和缩略语说明 }
        \cs_set:Npn \acknowledgmentsname { 致谢 }
        \cs_set:Npn \resumename
            { 个人简历、在学期间发表的学术论文及研究成果 }
        \cs_set:Npn \anonymousresumename
            { 在学期间发表的学术论文及研究成果 }
        \ctexset
            {
                contentsname   = 目录,
                figurename     = 图,
                tablename      = 表,
                listfigurename = 插图清单,
                listtablename  = 附表清单,
                appendixname   = 附录,
                bibname        = 参考文献,
            }
    }
%    \end{macrocode}
% \end{macro}
%
% \begin{macro}{\@@_lang_zh_mode_two:}
%
%    \begin{macrocode}
\cs_new_protected:Nn \@@_lang_zh_mode_two:
    {
        \ctexset
            {
                chapter/format = {
                    \mslinespread { 1 }
                    \heiti
                    \zihao { 3 }
                    \bfseries
                    \centering
                },
                section = {
                    name      = { 第, 节 },
                    number    = \chinese{section},
                    format    = {
                        \mslinespread { 1 }
                        \heiti
                        \zihao { 4 }
                        \bfseries
                        \centering
                    },
                },
                subsection = {
                    name      = { , 、 },
                    number    = \chinese{subsection},
                    indent    = \parindent,
                    aftername = { },
                    tocline   = {
                        \CTEXifname
                            { \protect { \csname CTEXthe##1 \endcsname } }
                            { }
                        ##2
                    },
                },
                subsubsection = {
                    name      = { （, ） },
                    number    = \chinese{subsubsection},
                    indent    = \parindent,
                    aftername = { },
                },
            }
    }
%    \end{macrocode}
% \end{macro}
%
% \begin{variable}{\c_@@_lang_zh_crefname_prop}
%
% 引用类型到中文名称的映射。
%
%    \begin{macrocode}
\prop_const_from_keyval:Nn \c_@@_lang_zh_crefname_prop
    {
        equation      = 式,
        figure        = 图,
        subfigure     = 子图,
        table         = 表,
        page          = 页,
        chapter       = 章,
        section       = 节,
        subsection    = 小节,
        subsubsection = 子节,
        appendix      = 附录,
        theorem       = 定理,
        lemma         = 引理,
        corollary     = 推论,
        proposition   = 命题,
        definition    = 定义,
        example       = 例,
        algorithm     = 算法,
        listing       = 列表,
        line          = 行,
    }
%    \end{macrocode}
% \end{variable}
%
% \begin{variable}{\c_@@_lang_zh_cref_counting_clist}
%
% 采用“第\meta{编号}\meta{量词}”表达方式的引用类型及其量词。
%
%    \begin{macrocode}
\clist_const:Nn \c_@@_lang_zh_cref_counting_clist
    {
        { page          } { 页   },
        { chapter       } { 章   },
        { section       } { 节   },
        { subsection    } { 小节 },
        { subsubsection } { 子节 },
    }
%    \end{macrocode}
% \end{variable}
%
% \begin{macro}{\@@_lang_zh_cref_set_counting:nn}
%
% 为一个“第\meta{编号}\meta{量词}”类型生成全部四组格式设置。
%
%    \begin{macrocode}
\cs_new_protected:Nn \@@_lang_zh_cref_set_counting:nn
    {
        \crefformat      {#1} { ##2第##1#2##3                }
        \crefrangeformat {#1} { ##3第##1#2##4至##5第##2#2##6 }
        \crefmultiformat {#1}
            { ##2第##1#2##3   } { 和##2第##1#2##3 }
            { ，##2第##1#2##3 } { 和##2第##1#2##3 }
        \crefrangemultiformat {#1}
            { ##3第##1#2##4至##5第##2#2##6   }
            { 和##3第##1#2##4至##5第##2#2##6 }
            { ，##3第##1#2##4至##5第##2#2##6 }
            { 和##3第##1#2##4至##5第##2#2##6 }
    }
%    \end{macrocode}
% \end{macro}
%
% \begin{macro}{\@@_lang_zh_cref:}
%
% 设置中文的交叉引用格式。
%
%    \begin{macrocode}
\cs_new_protected:Nn \@@_lang_zh_cref:
    {
        \labelformat { subfigure } { \thefigure（\thesubfigure） }
        \labelformat { equation  } { （\theequation）            }
        \prop_map_inline:Nn \c_@@_lang_zh_crefname_prop
            { \crefname {##1} {##2} {##2} }
        \clist_map_inline:Nn \c_@@_lang_zh_cref_counting_clist
            { \@@_lang_zh_cref_set_counting:nn ##1 }
%    \end{macrocode}
%
% 以下两个类型的量词位于编号之前，与“第\meta{编号}\meta{量词}”不同，因此单独设置。
%
%    \begin{macrocode}
        \crefformat { equation  } { 式##2##1##3 }
        \crefformat { subfigure } { 图##2##1##3 }
        \cs_set:Npn \crefpairconjunction        { ~和~ }
        \cs_set:Npn \crefmiddleconjunction      { ，   }
        \cs_set:Npn \creflastconjunction        { ~和~ }
        \cs_set:Npn \crefpairgroupconjunction   { ~和~ }
        \cs_set:Npn \crefmiddlegroupconjunction { ，   }
        \cs_set:Npn \creflastgroupconjunction   { ~和~ }
        \cs_set:Npn \crefrangeconjunction       { ~至~ }
    }
%    \end{macrocode}
% \end{macro}
